\documentclass[12pt]{article} 
\renewcommand{\thempfootnote}{\arabic{mpfootnote}}

\usepackage[utf8]{inputenc}

\usepackage{multicol}
\setlength{\columnsep}{1cm}

\usepackage{abc}
\usepackage{graphicx}
\usepackage{musicography}

\usepackage[margin=0.5in]{geometry} 
\geometry{letterpaper}

\usepackage{fancyhdr} 
\pagestyle{fancy} % options: empty , plain , fancy
\renewcommand{\headrulewidth}{0pt} % customise the layout...
\lhead{}\chead{}\rhead{}
\lfoot{}\cfoot{\thepage}\rfoot{}

\begin{document}

\thispagestyle{empty}

\begin{abc}[name=dryanddusty,options={-O= +c -F single.fmt}]
X:1
T:Dry and Dusty
K:D
|: AB  | "D"d2 d2 e2 de | f2 e2 dcBA | d2 d2 e2 de | "A"f2 a2 a2 AB | 
"D"d2 d2 e2 de | "(G)"f2 a2 f2-fe | "(A)"fedA "A"BAFB | "D"A2 D2 D2 :|
|: AB | "D"A2 F2 F2 AB | A2 D2 D2 AB | A2 F2 "A"E2 DF | "D"E2 D2 D2 AB | 
A2 F2 F2 AB | AFED "G"d2-dB | "D"A2 F2 "A"E2 DF | "D"E2 D2 D2 :|
\end{abc}
\vskip 0.5in
\noindent\makebox[\textwidth][c]{%
  \begin{minipage}{7.0in}
    \setlength{\parindent}{1.5em}
    
    Dry and Dusty is a popular Old Time tune said to be from the
    Arkansas/Missouri Ozark area. This tune is often played in
    alternative fiddle tunings such as AEae (Arkansas fiddler Lon
    Jordan playing in A), DDad (Eck Robertson, Benny Thomasson, the
    Morrison twins - Apsie Sherdon Morrison and Abbie Sherman
    Morrison), ADad (Benny Thomasson, Morrison Twin Brothers String
    Band), as well as the usual
    GDae.\footnote{https://tunearch.org/wiki/Annotation:Dry\_and\_Dusty\_(1)}
    At one point the tune was so well known in the Ozark region that
    Ken Perlman, noted folklorist and clawhammer banjoist, claimed
    that the tune was played whenever the fiddler wanted tell the
    staff that they were ``Dry and Dusty'' and wanted a drink.

    The version transcribed is a blend of a few different versions, as
    tunes learned in jams or at festivals often are.

  \end{minipage}
}


\end{document}
